\answer{ex:20:1}{$t_{\text{醒}}=\tau_r\ln\frac{1-L}{1-H}=16.8$~h，$t_{\text{睡}}=\tau_d\ln\frac HL=5.8$~h，周期 $22.5$~h；阈值随昼夜的摆动把振荡器拉到 $24$~h——这就是锁相~\eqref{eq:pll}。}
\answer{ex:20:2}{(a)~$L=0.111$；若 $L=0.092$，睡眠会持续 $9.6$~h。(b)~增长 $22\%$（$1/0.82=1.22$），而不是 $18\%$。}
\answer{ex:20:3}{$H=\frac{p-1}{p-1/q}=0.623$，$L=H/q=0.093$，其中 $p=e^{16/\tau_r}$，$q=e^{8/\tau_d}$。在 $4$~h 后被叫醒，相位会永久偏移（入睡提前 $7.2$~h）；阈值摆动时，相位在两三天内恢复。}
\answer{ex:20:4}{(a)~$r_\pm=\frac12\bigl(1\pm\sqrt{1-4k/w}\bigr)=0.947,\ 0.053$；$a_{\text{上}}=3.51$，$a_{\text{下}}=0.49$。(b)~工作 $\approx0.33$~s，静默 $\approx0.59$~s，周期 $\approx0.93$~s，工作比例 $0.36$。}
\answer{ex:20:5}{$a_{\text{上}}/r_+<b<a_{\text{下}}/r_-$：$3.71<b<9.20$。\nt{ACET} 使 $b$ 降到 $3.71$ 以下：交点移到上分支，工作状态（$r\approx1$）变得稳定。}
\answer{ex:20:6}{$D=24$、$32$、$40$、$48$、$64$ 时分别为 $4.9$、$6.8$、$8.8$、$5.5$、$8.9$~h：时长由入睡时的昼夜相位决定，而不是由欠债决定。}
\answer{ex:20:7}{学完 B 后，A 上的误差平均为 $0.51$（$20$ 组数据中从 $0.19$ 到 $1.08$；零输出的误差为 $1$）。交错学习时两个误差都小于 $10^{-4}$。}
\answer{ex:20:8}{开放问题。均匀缩放保留权重的比值，但只有当阈值按同一比例缩放时，才能保留阈值神经元的回答；交错重放则改变比值。大接触保持不变，这与纯粹的均匀缩放相矛盾。}
